\scanpage{91}\mspage{T-5}

\uline{Cas}\quad $j=0$\quad (i.e. $a_{4}=0$) .

Ici $C\approx C'$ équivaut à $j=j'$ et $a_{6}/a'_{6}\in k^{\ast 6}$ .

Les automorphismes sont donnés par les $u$ tels que $u^{6}=1$.

Courbe type : $y^{2}=x^{3}-1$.

4. \uline{Forme canonique} ( $p=3$)

Du fait que $2$ est inversible, on peut tout d'abord écrire l'équation sous la forme :
\[
 y^{2}=x^{3}+a_{2}x^{2}+a_{4}x+a_{6}=f(x).
\]

On a :
\[
\begin{aligned}
 &b_{2}=a_{2}\ ,\quad b_{4}=-a_{4}\ ,\quad b_{8}=a_{4}^{2}-a_{2}a_{6}
 \qquad\quad \pi=-dx/y\\
 &c_{4}=a_{2}^{2}\ ,\quad c_{6}=a_{2}^{3}\ ,\quad
 \handwritten{\Delta}=a_{2}^{2}a_{4}^{2}-a_{2}^{3}a_{6}-a_{4}^{3}\ .
\end{aligned}
\]

Ici encore, $\handwritten{\Delta}$ est, à un facteur inversible près, le discriminant de $f$. Le fait que $C$ soit elliptique se traduit donc par $\handwritten{\Delta}\ne 0$, et $j=a_{2}^{6}/\handwritten{\Delta}$ . Deux courbes $C$ et $C'$ de même $j$ sont isomorphes après extension séparable de $k$. Plus précisément:

\uline{Cas général} : $j\ne 0$\quad (i.e. $a_{2}\ne 0$).

On peut faire disparaître le terme en $x$. Cela donne la forme réduite :
\[
 y^{2}=x^{3}+a_{2}x^{2}+a_{6}\ ,\quad
 \handwritten{\Delta}=-a_{2}^{3}a_{6}\ ,\quad j=-a_{2}^{3}/a_{6}\ .
\]

Un isomorphisme $f:C'\mathrel{\handwritten{\to}}C$ est donné par
\[
 x\;f=u^{2}x'\qquad y\unsure{\mathord{.}|\circ|\phantom{\circ}}f=u^{3}x'\ .
\]

Pour que $C\approx C'$ il faut et il suffit que $j=j'$ et que $a_{2}/a'_{2}\in k^{\ast 2}$ . Les automorphismes de $C$ correspondent à $u=\pm 1$.
% archive sheet number: 90

